\documentclass[11pt,a4paper]{article}

\usepackage{graphicx}
\usepackage{amsmath}
\usepackage{booktabs}
\usepackage{hyperref}
\usepackage{geometry}

\geometry{margin=1in}


\title{
From LeetCode to Prompt Engineering:
Investigating the Impact of LLM-Assisted Programming on Algorithmic Problem Solving
}


\author{
Chunru Qiu\\
Department of Computer Science\\
Concordia University
}


\date{2026}


\begin{document}


\maketitle


\begin{abstract}

Large Language Models (LLMs) are rapidly transforming software engineering
practices by assisting developers in code generation, debugging, and optimization.
Traditionally, computer science students prepare for software engineering
interviews through extensive algorithm and data structure practice on platforms
such as LeetCode.

This research investigates whether optimized prompt engineering can improve
algorithmic problem-solving performance and potentially change the traditional
algorithm training paradigm.

We propose an experimental framework comparing different AI-assisted
problem-solving approaches, including baseline prompting, optimized algorithm
reasoning prompts, and structured AI-assisted workflows.

The study evaluates generated solutions based on correctness, algorithm
selection, computational complexity, and runtime performance.

\end{abstract}



\section{Introduction}


Software engineering has historically emphasized strong foundations in
data structures, algorithms, and computational complexity.

Platforms such as LeetCode have become widely used by computer science students
to develop algorithmic problem-solving skills and prepare for technical interviews.


The traditional learning process can be summarized as:

\begin{center}
Data Structures
$\rightarrow$
Algorithms
$\rightarrow$
Problem Practice
$\rightarrow$
Technical Interview Preparation
\end{center}


However, recent advances in Large Language Models have introduced a new
software development workflow.

Developers can now interact with AI systems to generate solutions, analyze
complexity, and optimize implementations through carefully designed prompts.


This creates an important research question:

\begin{quote}
Can prompt engineering become an alternative or complementary skill to
traditional algorithm practice?
\end{quote}


This research explores how AI-assisted programming may influence future
software engineering education and interview preparation.



\section{Research Question}


The primary research question is:


\textbf{
Can optimized prompt engineering improve algorithmic problem-solving
performance compared with traditional LeetCode-style approaches?
}


The study investigates three sub-questions:


\begin{enumerate}

\item
Can optimized prompts improve the correctness of generated solutions?


\item
Can prompt engineering improve algorithm and data structure selection?


\item
How may AI-assisted programming change the skills required by future
software engineers?

\end{enumerate}



\section{Background}


\subsection{Traditional Algorithm Training}


Traditional computer science education emphasizes:

\begin{itemize}
\item Data structures
\item Algorithm design
\item Complexity analysis
\item Independent implementation
\end{itemize}



\subsection{LLM-Assisted Programming}


Large Language Models introduce a different development paradigm:


\begin{center}

Problem Understanding

$\rightarrow$

Prompt Engineering

$\rightarrow$

AI Generated Solution

$\rightarrow$

Human Verification

$\rightarrow$

Optimization

\end{center}

\section{Related Work}

\subsection{Algorithmic Problem Solving Education}

Algorithmic problem solving has been considered a fundamental skill in
computer science education and software engineering recruitment.

Traditional approaches emphasize the understanding of data structures,
algorithm design techniques, and computational complexity analysis.
Platforms such as LeetCode have become widely adopted tools for students
to practice programming problems and prepare for technical interviews.

Previous educational research has investigated how repeated programming
practice improves algorithmic thinking, pattern recognition, and problem
solving ability.

However, traditional algorithm training requires significant time investment,
as students often solve hundreds of problems to develop sufficient experience.


\subsection{Large Language Models for Code Generation}

Recent advances in Large Language Models (LLMs) have significantly changed
software development workflows.

Models such as GPT-based systems and other code-capable LLMs have demonstrated
the ability to generate code, explain algorithms, and assist developers in
debugging tasks.

Previous studies have evaluated LLM performance on programming benchmarks and
found that these models can solve many software engineering tasks while still
facing challenges in correctness, reasoning, and complex algorithm selection.

These limitations suggest that effective interaction strategies between humans
and LLMs remain an important research direction.


\subsection{Prompt Engineering for Software Development}

Prompt engineering has emerged as a method to improve LLM performance by
providing structured instructions and reasoning guidance.

Research has shown that carefully designed prompts can improve model outputs
in tasks involving reasoning, planning, and code generation.

In software engineering contexts, prompts that include requirements analysis,
algorithm planning, and verification steps may help LLMs produce higher-quality
solutions.

However, limited research has investigated whether prompt engineering can
change the role of traditional algorithm training, especially in the context
of software engineering interview preparation.


\subsection{Research Gap}

Existing studies have primarily focused on evaluating whether LLMs can generate
correct code or improve developer productivity.

However, fewer studies have examined the broader question:

\begin{quote}
Will algorithmic training practices for future software engineers shift from
traditional problem solving toward AI-assisted prompt engineering workflows?
\end{quote}


This research addresses this gap by comparing different prompting strategies
for algorithmic problem solving and investigating how AI-assisted programming
may influence future software engineering skill requirements.



\section{Methodology}


This research proposes a benchmark framework to evaluate different prompting
strategies for algorithmic problem solving.


The experiment compares three approaches:


\subsection{Baseline Prompt}


A minimal prompt asking the model to directly solve the problem.


\subsection{Optimized Prompt}


A structured prompt requiring:

\begin{itemize}
\item Constraint analysis
\item Algorithm pattern recognition
\item Data structure selection
\item Complexity comparison
\end{itemize}



\subsection{AI-Assisted Engineering Workflow}


A multi-step workflow where the model:

\begin{itemize}
\item Generates multiple approaches
\item Compares trade-offs
\item Verifies edge cases
\item Produces optimized solutions
\end{itemize}


\section{Evaluation Methodology}


To evaluate the impact of prompt engineering on algorithmic problem solving,
this study evaluates generated solutions using multiple quantitative metrics.

The evaluation focuses on solution correctness, algorithm quality, computational
efficiency, and execution performance.



\subsection{Correctness Evaluation}


The primary evaluation metric is solution correctness.


Each generated Python solution is executed against predefined test cases from
the benchmark dataset.


The correctness score is calculated as:


\[
Correctness = \frac{Passed\ Test\ Cases}{Total\ Test\ Cases}
\]


A higher correctness score indicates that the generated solution successfully
solves more programming tasks.



\subsection{Algorithm Selection Evaluation}


Beyond correctness, this research evaluates whether the generated solution
selects an appropriate algorithmic approach.


The evaluation compares:

\begin{itemize}

\item Selected algorithm pattern.

\item Expected optimal algorithm.

\item Data structure usage.

\end{itemize}


For example, a solution using a hash table for a two-sum problem is considered
more appropriate than a brute-force nested loop solution.



\subsection{Complexity Evaluation}


Generated solutions are analyzed based on computational complexity.


The evaluation considers:


\begin{itemize}

\item Time complexity.

\item Space complexity.

\end{itemize}


The generated complexity is compared with the expected optimal complexity
provided by the benchmark dataset.


A complexity score is assigned:


\[
Complexity\ Score =
\frac{Generated\ Efficiency}{Optimal\ Efficiency}
\]


This metric evaluates whether prompt engineering helps models produce more
optimized algorithms.



\subsection{Runtime Performance}


Each generated program is executed in a controlled environment.


Runtime performance is measured using:

\begin{itemize}

\item Execution time.

\item Memory usage.

\item Successful completion.

\end{itemize}


Runtime measurements are averaged across multiple test executions to reduce
random variation.



\subsection{Prompt Strategy Comparison}


The experiment compares three prompting strategies:


\begin{table}[h]

\centering

\begin{tabular}{lll}

\toprule

Strategy & Description & Purpose \\

\midrule

Baseline & Direct problem solving prompt & Measure default LLM capability \\

Optimized & Structured algorithm reasoning prompt & Evaluate prompt improvement \\

AI-Assisted & Multi-step engineering workflow & Evaluate collaborative workflow \\

\bottomrule

\end{tabular}


\caption{Prompt evaluation strategies}

\end{table}



The final comparison analyzes whether more structured prompting strategies
provide measurable improvements over direct LLM interaction.



\subsection{Statistical Analysis}


The experiment reports:

\begin{itemize}

\item Average correctness score.

\item Average runtime.

\item Complexity improvement.

\item Performance differences among prompt strategies.

\end{itemize}


These measurements are used to determine whether prompt engineering provides
significant benefits for algorithmic problem solving.

\section{Experimental Design}


The benchmark dataset will contain LeetCode-style problems from multiple
algorithm categories:


\begin{itemize}
\item Arrays
\item Hash Tables
\item Linked Lists
\item Trees
\item Graphs
\item Dynamic Programming
\end{itemize}



Each generated solution will be evaluated using:


\begin{itemize}

\item Correctness
\item Time complexity
\item Space complexity
\item Runtime performance
\item Data structure selection

\end{itemize}



\section{Expected Contribution}


This research aims to provide:


\begin{enumerate}

\item
A benchmark framework for evaluating LLM-assisted algorithm problem solving.


\item
An analysis of how prompt engineering affects algorithm design decisions.


\item
Insights into how software engineering skills may evolve in the LLM era.

\end{enumerate}



\section{Future Work}


Future work includes:


\begin{itemize}

\item Implementing an automated evaluation pipeline.

\item Testing multiple Large Language Models.

\item Expanding the benchmark dataset.

\item Analyzing experimental results.

\item Publishing findings as an open research report.

\end{itemize}



\section{Conclusion}


Large Language Models are changing how developers interact with programming
tasks.

This research investigates whether prompt engineering can become an important
skill in future software engineering workflows while examining its relationship
with traditional algorithm training.

\bibliographystyle{plain}
\bibliography{references}

\end{document}
